
\documentclass[12pt,a4paper]{article}
\usepackage[spanish]{babel}
\usepackage[utf8]{inputenc}
\usepackage[T1]{fontenc}
\usepackage{lmodern}
\usepackage{geometry}
\geometry{margin=1in}
\usepackage{hyperref}
\usepackage{enumitem}
\usepackage{graphicx}
\usepackage{xcolor}
\usepackage{listings}
\usepackage{caption}
\usepackage{float}

\hypersetup{
  colorlinks=true,
  linkcolor=blue,
  urlcolor=blue,
  citecolor=blue
}

\lstdefinestyle{mypy}{
  language=Python,
  basicstyle=\ttfamily\small,
  keywordstyle=\color{blue}\bfseries,
  stringstyle=\color{teal!70!black},
  commentstyle=\color{gray!70},
  showstringspaces=false,
  frame=single,
  breaklines=true,
  tabsize=2
}

\title{\textbf{Reto (sin solución): Fashion-MNIST desde cero con NumPy}}
\author{Curso de Aprendizaje Automático}
\date{\today}

\begin{document}
\maketitle

\section*{Objetivo}
Implementar \textbf{desde cero (solo NumPy)} un clasificador para \textbf{Fashion-MNIST} (10 clases de prendas). No se permite usar librerías de Deep Learning (Keras/TensorFlow/PyTorch/JAX). El reto incluye carga de datos, preprocesamiento, construcción de una \textbf{red neuronal multicapa}, entrenamiento con \textbf{mini-batches}, visualizaciones y análisis de errores.

\section*{Reglas y Alcance}
\begin{itemize}[leftmargin=*,nosep]
  \item \textbf{Permitido}: \texttt{numpy}, \texttt{pandas}, \texttt{matplotlib}. (\texttt{sklearn} solo para métricas como matriz de confusión).
  \item \textbf{Prohibido}: Keras, TensorFlow, PyTorch, JAX u otras librerías de DL.
  \item \textbf{Datos}: Fashion-MNIST (formato MNIST: 28$\times$28 gris, 10 clases). Use la versión CSV (columna \texttt{label} + 784 píxeles) o convierta desde \texttt{idx}/\texttt{npz}.
  \item \textbf{Entrega}: código reproducible + figuras + breve informe (README).
\end{itemize}

\section*{Clases (para títulos en gráficas)}
\begin{center}
0=T-shirt/top, 1=Trouser, 2=Pullover, 3=Dress, 4=Coat, 5=Sandal, 6=Shirt, 7=Sneaker, 8=Bag, 9=Ankle boot
\end{center}

\section*{Estructura sugerida de proyecto}
\begin{verbatim}
fashion_reto/
├─ data/                  # train.csv (o .npz/.idx)
├─ figs/                  # imágenes de salida
├─ src/
│  ├─ reto_fashion.py     # implementación principal
│  └─ utils.py            # helpers (opcional)
└─ README.md              # cómo ejecutar y hallazgos
\end{verbatim}

\section*{Entregables}
Guarde las figuras en \texttt{figs/} con estos nombres:
\begin{itemize}[leftmargin=*,nosep]
  \item \texttt{fig\_fm\_imagen\_0\_con\_valores.png}
  \item \texttt{fig\_fm\_10\_ejemplos.png}
  \item \texttt{fig\_fm\_loss\_train\_val.png}
  \item \texttt{fig\_fm\_confusion\_matrix.png}
  \item \texttt{fig\_fm\_errores\_grid.png}
\end{itemize}
Reporte en README: accuracy en \textit{train/valid/test} y 2--3 observaciones sobre confusiones típicas.

\section*{Pasos (qué hacer, \underline{sin solución})}

\subsection*{0) Preparación}
\begin{enumerate}[leftmargin=*,nosep,label=\arabic*)]
  \item Cree un entorno e instale: \texttt{numpy pandas matplotlib}.
  \item Fije semilla: \texttt{np.random.seed(100)}.
\end{enumerate}

\subsection*{1) Carga + EDA breve}
\begin{enumerate}[leftmargin=*,nosep,label=\alph*)]
  \item Lea \texttt{data/train.csv}. Separe: \texttt{y = df['label'].values}, \texttt{X\_raw = df.drop('label', axis=1).values}.
  \item Imprima \texttt{X\_raw.shape}, distribución de clases.
  \item Visualice:
    \begin{itemize}[leftmargin=*,nosep]
      \item Imagen 0 (28$\times$28) con \textbf{valores por píxel} (0--255) sobre la celda.
      \item Cuadrícula 2$\times$5 con los \textbf{primeros 10} ejemplos (título: nombre de clase).
    \end{itemize}
\end{enumerate}

\subsection*{2) Preprocesamiento}
\begin{enumerate}[leftmargin=*,nosep,label=\alph*)]
  \item Normalice: \texttt{X = X\_raw.astype(np.float32)/255.0}.
  \item Codifique etiquetas con \textbf{One-Hot}: \texttt{Y = one\_hot(y, num\_classes=10)}.
  \item Divida en \textbf{train/valid/test} (80/10/10). Si no usa \texttt{sklearn}, baraje índices manualmente.
\end{enumerate}

\subsection*{3) Arquitectura}
Use base: \textbf{784 $\rightarrow$ 64 $\rightarrow$ 32 $\rightarrow$ 10}. Active en ocultas con \textbf{ReLU} o \textbf{Sigmoid}. En salida, \textbf{Softmax} (recomendado) y pérdida \textbf{Cross-Entropy}.

\subsection*{4) Inicialización}
Inicialice pesos pequeños (escala 0.01) y sesgos en 0. Imprima número de parámetros por capa.

\subsection*{5) Forward}
Implemente \texttt{forward(X, params)} que devuelva:
\begin{itemize}[leftmargin=*,nosep]
  \item \texttt{P}: probabilidades de clase por lote.
  \item \texttt{cache}: diccionario con intermedios (\texttt{a0,z0,a1,z1,a2,z2}) para backprop.
\end{itemize}

\subsection*{6) Pérdida}
Implemente \textbf{Cross-Entropy} multiclase con estabilidad numérica (restar \texttt{max} por fila antes de \texttt{exp}).

\subsection*{7) Backpropagation}
Calcule gradientes para todas las capas. Con Softmax+CE:
\begin{center}
$\delta^{(L)} = \dfrac{P - Y}{B}$, \quad
$\delta^{(1)} = (\delta^{(L)} W_2^\top) \odot g'(z^{(1)})$, \quad
$\delta^{(0)} = (\delta^{(1)} W_1^\top) \odot g'(z^{(0)})$
\end{center}
donde $g$ es la activación de las ocultas y $B$ el tamaño de \textit{batch}.

\subsection*{8) Entrenamiento (mini-batch)}
Baraje cada época. Para cada mini-batch: \texttt{forward $\rightarrow$ loss $\rightarrow$ backward $\rightarrow$ update}.
Registre \texttt{loss\_train} y \texttt{loss\_val}. Grafique y guarde \texttt{fig\_fm\_loss\_train\_val.png}.

\subsection*{9) Evaluación + Errores}
\begin{itemize}[leftmargin=*,nosep]
  \item Calcule \textbf{accuracy} en train/valid/test.
  \item Genere \textbf{matriz de confusión} (\texttt{sklearn.metrics.confusion\_matrix} permitido).
  \item Muestre 12 errores en cuadrícula 3$\times$4 con títulos: \texttt{Real=\textless nombre\textgreater, Pred=\textless nombre\textgreater}. Guarde \texttt{fig\_fm\_errores\_grid.png}.
\end{itemize}

\subsection*{10) Experimentos (elija 2+)}
Cambie tamaño de capas (p.ej., 128/64), compare \textbf{ReLU vs Sigmoid}, pruebe distintas tasas de aprendizaje, agregue \textbf{L2} o \textbf{early stopping}. Documente hallazgos en README.

\section*{Plantilla (esqueleto, sin solución)}
\begin{lstlisting}[style=mypy,caption={Esqueleto de src/reto_fashion.py con TODOs}]
import numpy as np
import pandas as pd
import matplotlib.pyplot as plt

CLASS_NAMES = ["T-shirt/top","Trouser","Pullover","Dress","Coat",
               "Sandal","Shirt","Sneaker","Bag","Ankle boot"]

# ---------- utils ----------
def one_hot(y, num_classes=10):
    # TODO: devolver matriz (N,num_classes) one-hot.
    raise NotImplementedError

def relu(z):
    # TODO: implementar ReLU.
    raise NotImplementedError

def relu_derivative(z):
    # TODO: derivada de ReLU.
    raise NotImplementedError

def softmax(z):
    # TODO: softmax estable numéricamente.
    raise NotImplementedError

def cross_entropy(y_true, y_pred):
    # TODO: pérdida CE multiclase estable.
    raise NotImplementedError

# ---------- modelo ----------
def init_params(input_dim=784, h1=64, h2=32, num_classes=10, scale=0.01):
    # TODO: inicializar W0,W1,W2 y b0,b1,b2.
    raise NotImplementedError

def forward(X, params, activation="relu"):
    # TODO: calcular a0,z0,a1,z1,a2,z2 y salida P.
    # Usa 'activation' para ocultas; salida softmax.
    # Retorna P y cache.
    raise NotImplementedError

def backward(Y, cache, params, activation="relu"):
    # TODO: gradientes con softmax+CE:
    #   delta2 = (P - Y) / B
    #   delta1 = (delta2 @ W2.T) * g'(z1)
    #   delta0 = (delta1 @ W1.T) * g'(z0)
    raise NotImplementedError

def update(params, grads, lr):
    # TODO: descenso por gradiente.
    raise NotImplementedError

def accuracy(y_true_labels, y_pred_probs):
    # TODO: argmax y promedio de aciertos.
    raise NotImplementedError

# ---------- visualizaciones obligatorias ----------
def show_first_image_with_values(X_raw, y, outpath="figs/fig_fm_imagen_0_con_valores.png"):
    # TODO: imagen 0 con valores 0..255 impresos en cada pixel.
    raise NotImplementedError

def show_first_10_grid(X_raw, y, outpath="figs/fig_fm_10_ejemplos.png"):
    # TODO: cuadrícula 2x5 con títulos CLASS_NAMES[y[i]].
    raise NotImplementedError

def show_errors_grid(X_raw, y_true, y_pred, k=12, outpath="figs/fig_fm_errores_grid.png"):
    # TODO: 3x4 ejemplos erróneos con títulos 'Real= , Pred='.
    raise NotImplementedError

# ---------- entrenamiento ----------
def train(X_train, Y_train, X_val, Y_val, params, lr=0.01, batch_size=128, epochs=20, activation="relu"):
    # TODO: bucle mini-batch:
    #   - barajar por época
    #   - forward -> loss -> backward -> update
    #   - registrar pérdidas y devolver history
    raise NotImplementedError

def main():
    np.random.seed(100)
    df = pd.read_csv("data/train.csv")
    y = df["label"].values
    X_raw = df.drop("label", axis=1).values.astype(np.float32)

    show_first_image_with_values(X_raw, y)
    show_first_10_grid(X_raw, y)

    X = X_raw / 255.0
    Y = one_hot(y, num_classes=10)

    # TODO: split train/val/test
    # TODO: params = init_params()
    # TODO: history = train(...)
    # TODO: evaluación + matriz de confusión + show_errors_grid(...)
    # TODO: guardar curva de pérdidas

if __name__ == "__main__":
    main()
\end{lstlisting}

\section*{Criterios de evaluación (sugerencia)}
\begin{itemize}[leftmargin=*,nosep]
  \item \textbf{Correctitud técnica (40\%)}: forward/backprop correctos, convergencia, métricas.
  \item \textbf{Calidad de código (20\%)}: claridad, modularidad, comentarios, reproducibilidad.
  \item \textbf{Visualizaciones (20\%)}: figuras solicitadas, legibles y guardadas.
  \item \textbf{Análisis (20\%)}: interpretación de errores y propuestas de mejora.
\end{itemize}

\section*{Cómo compilar}
Ejecute: \texttt{pdflatex reto\_fashion\_mnist.tex} (dos pasadas recomendadas).

\bigskip
\noindent\textbf{Nota}: Este documento describe el \emph{reto} y proporciona una plantilla con \texttt{TODOs}. No incluye soluciones.

\end{document}
